\documentclass[10pt]{article}
\usepackage[UTF8]{ctex}
\usepackage{amsmath, amssymb}
\usepackage{geometry}
\geometry{a4paper, margin=1.5cm}
\usepackage{enumitem}

\title{计数原理习题与解答}
\author{}
\date{}

\begin{document}
\maketitle

\section*{习题部分}

\subsection*{题1（数字组数）}
用数字 0, 1, 2, 3, 4, 5 组成无重复数字的四位数，其中偶数的个数是多少？

\subsection*{题2（排队问题）}
6 个人站成一排，求下列情况的排法数：
\begin{enumerate}
\item 甲、乙、丙三人必须站在一起；
\item 甲不站在排头，乙不站在排尾。
\end{enumerate}

\subsection*{题3（分配问题）}
将 5 本不同的书分给甲、乙、丙三人，每人至少一本，有多少种不同的分法？

\subsection*{题4（涂色问题）}
用 4 种不同颜色给如图所示的 5 个区域涂色，要求相邻区域颜色不同。区域编号为 1-2-3-4-5 依次相邻（直线型），且区域 1 与区域 5 不相邻。求不同的涂色方案数。

\subsection*{题5（至多至少问题）}
某班有 30 名学生，其中 10 名女生，20 名男生。现要选出 5 人组成代表队，要求至少有 2 名女生，有多少种选法？

\subsection*{题6（分组分配，均匀分组）}
将 4 名医生和 4 名护士平均分成 4 组，每组 1 名医生和 1 名护士，分别派往 4 所不同的学校，有多少种分配方案？

\subsection*{题7（隔板法）}
求方程 \(x_1 + x_2 + x_3 + x_4 = 12\) 的非负整数解的个数。

\subsection*{题8（综合，分类与分步）}
从 5 双不同尺码的鞋子中任取 4 只，求下列事件的取法数：
\begin{enumerate}
\item 恰好有 2 只成双（即一双完整，另外两只不成双）；
\item 4 只鞋中至少有 2 只成双。
\end{enumerate}

\section*{解答部分}

\subsection*{题1 解答}
\textbf{思路}：四位偶数，个位必须是偶数。数字含 0，需分个位是否为 0 两类。

\textbf{解}：
\begin{itemize}
\item \textbf{第一类}：个位为 0。
  千位从 1~5 中选（5 种），百位从剩余 4 个中选（4 种），十位从剩余 3 个中选（3 种）。
  方法数：\(5 \times 4 \times 3 = 60\)。
\item \textbf{第二类}：个位为 2 或 4。
  个位有 2 种选择。千位不能为 0 且不能与个位相同，故从剩下的 4 个非 0 且非个位数字中选（4 种）。百位从剩余 4 个数字中选（含 0，4 种），十位从剩余 3 个中选（3 种）。
  方法数：\(2 \times 4 \times 4 \times 3 = 96\)。
\end{itemize}
总数：\(60 + 96 = 156\)。

\textbf{答}：156 个。

\subsection*{题2 解答}
\textbf{(1) 甲、乙、丙三人必须站在一起}

\textbf{解}：捆绑法。将甲、乙、丙看作一个整体，内部有 \(3! = 6\) 种排法。
整体与其余 3 人共 4 个元素，全排列 \(4! = 24\)。
总数：\(6 \times 24 = 144\)。

\textbf{(2) 甲不站排头，乙不站排尾}

\textbf{解}：间接法。总排法 \(6! = 720\)。
\begin{itemize}
\item 甲在排头：固定甲在第一位，其余 5 人全排列 \(5! = 120\)。
\item 乙在排尾：固定乙在最后一位，其余 5 人全排列 \(5! = 120\)。
\item 甲在排头且乙在排尾：固定甲在首位、乙在末位，中间 4 人全排列 \(4! = 24\)。
\end{itemize}
由容斥原理，不符合条件的方法数为 \(120 + 120 - 24 = 216\)。
符合条件的方法数：\(720 - 216 = 504\)。

\textbf{答}：(1) 144 种；(2) 504 种。

\subsection*{题3 解答}
\textbf{思路}：5 本不同书分给 3 个不同人，每人至少 1 本 → 先分组（将 5 本书分成 3 组，组的大小为 3,1,1 或 2,2,1），再分配。

\textbf{解}：
\begin{itemize}
\item \textbf{情形一}：3,1,1 分组
  选 3 本书为一组：\(C_5^3 = 10\)，剩下两本书各自成组（自然区分，因为书不同，且两单本组不同）。分组数为 10。
  将三组（有区别，因书不同）分配给 3 个人：\(3! = 6\)。
  小计：\(10 \times 6 = 60\)。
\item \textbf{情形二}：2,2,1 分组
  先选 1 本书作为单本组：\(C_5^1 = 5\)。
  剩下 4 本书分成两堆各 2 本：注意均匀分组要除 \(2!\)，即 \(\frac{C_4^2 C_2^2}{2!} = \frac{6 \times 1}{2} = 3\)。
  故分组数为 \(5 \times 3 = 15\)。
  分配给 3 个人：\(3! = 6\)。
  小计：\(15 \times 6 = 90\)。
\end{itemize}
总数：\(60 + 90 = 150\)。

\textbf{答}：150 种。

\subsection*{题4 解答}
\textbf{思路}：直线型 5 个区域，依次相邻（1-2,2-3,3-4,4-5 相邻，1与5不相邻）。用分步涂色，注意区域 3 与 1 不相邻，但与 2,4 相邻。

\textbf{解}：
\begin{itemize}
\item 区域 1：4 种颜色。
\item 区域 2：不能与 1 同色，有 3 种。
\item 区域 3：不能与 2 同色，可以与 1 同色，故有 3 种（除去区域 2 的颜色）。
\item 区域 4：不能与 3 同色，有 3 种。
\item 区域 5：不能与 4 同色（但与 1 不相邻，故不限制与 1 同色），有 3 种。
\end{itemize}
总数：\(4 \times 3 \times 3 \times 3 \times 3 = 4 \times 3^4 = 4 \times 81 = 324\)。

\textbf{答}：324 种。

\subsection*{题5 解答}
\textbf{思路}：从 10 女 20 男中选 5 人，至少 2 女。可用直接分类或间接法。

\textbf{解}（直接分类）：
\begin{itemize}
\item 2 女 3 男：\(C_{10}^2 \cdot C_{20}^3\)
\item 3 女 2 男：\(C_{10}^3 \cdot C_{20}^2\)
\item 4 女 1 男：\(C_{10}^4 \cdot C_{20}^1\)
\item 5 女 0 男：\(C_{10}^5\)
\end{itemize}
计算：
\[
C_{10}^2 = 45,\ C_{20}^3 = 1140 \quad \Rightarrow \quad 45 \times 1140 = 51300
\]
\[
C_{10}^3 = 120,\ C_{20}^2 = 190 \quad \Rightarrow \quad 120 \times 190 = 22800
\]
\[
C_{10}^4 = 210,\ C_{20}^1 = 20 \quad \Rightarrow \quad 210 \times 20 = 4200
\]
\[
C_{10}^5 = 252
\]
总和：\(51300 + 22800 = 74100\)，再加 \(4200 = 78300\)，再加 \(252 = 78552\)。

\textbf{答}：78552 种。

\subsection*{题6 解答}
\textbf{思路}：4 名医生、4 名护士，每组 1 医 1 护，派往 4 所不同学校。可理解为：先分配医生（全排列），再分配护士（全排列），或者直接分步。

\textbf{解}：
\begin{itemize}
\item 将 4 名医生分配到 4 所学校，每校 1 人：\(4! = 24\) 种。
\item 再将 4 名护士分配到 4 所学校，每校 1 人：\(4! = 24\) 种。
\end{itemize}
由乘法原理，总方案数：\(24 \times 24 = 576\)。

\textbf{答}：576 种。

\subsection*{题7 解答}
\textbf{思路}：方程 \(x_1+x_2+x_3+x_4=12\)，\(x_i \ge 0\) 整数解，用隔板法（相同元素分配问题）。
将 12 个相同的“1”分成 4 组（允许空组）。等价于 12 个球和 3 块隔板的全排列中，隔板位置决定分组。
标准公式：解数为 \(C_{12+4-1}^{4-1} = C_{15}^3\)。

计算：
\[
C_{15}^3 = \frac{15\times14\times13}{3\times2\times1} = 455.
\]

\textbf{答}：455 个解。

\subsection*{题8 解答}
\textbf{思路}：5 双不同尺码的鞋子，每双 2 只（左右脚对称，但通常视为不同个体）。 “成双”意味着取到同一双的两只。

\textbf{解}：
\begin{enumerate}
\item \textbf{恰好有 2 只成双（即一双完整，另两只不成双）}
  \begin{itemize}
  \item 第一步：选哪一双成双：\(C_5^1 = 5\) 种。
  \item 第二步：从剩下的 4 双中选 2 双，再从每双中选 1 只（不能选整双，否则就有两双了）。
    选 2 双：\(C_4^2 = 6\) 种。
    在选出的每双中，有 2 种选法（左或右），故 \(2^2 = 4\) 种。
    第二步共 \(6 \times 4 = 24\) 种。
  \item 总数：\(5 \times 24 = 120\)。
  \end{itemize}
\item \textbf{至少有 2 只成双}
  包含：恰好 1 双（即上问结果），恰好 2 双。
  \begin{itemize}
  \item 恰好 2 双：选 2 双，每双全取：\(C_5^2 = 10\)。
  \item 恰好 1 双：120。
  \end{itemize}
  总数 = \(120 + 10 = 130\)。

  \textbf{验证}：总取法 \(C_{10}^4 = 210\)。不成对（4 只来自不同双）：从 5 双中选 4 双，每双中选 1 只：\(C_5^4 \times 2^4 = 5 \times 16 = 80\)。则至少有一双：\(210 - 80 = 130\)。一致。
\end{enumerate}

\textbf{答}：(1) 120 种；(2) 130 种。

\section*{结束语}
以上习题覆盖了计数原理的常见高考题型。建议每做完一题，对照解答中的思路，反思自己是否准确识别了“分类/分步”“有序/无序”“元素是否相同”等关键点。通过适量练习，逐步形成稳定的解题思维框架。

\end{document}